% =====================================================================
%  RAPPORT DE PROJET — E-STORE
%  Module : Full-Stack
%  Encadrant : Pr. Omar Zahour
%  Établissement : FSBM — Université Hassan II de Casablanca
%  Auteurs : Akram Belmoussa & Nouhaila Ben Soumane
%  Année universitaire : 2025-2026
% =====================================================================
%
%  Compilation :
%    pdflatex rapport-estore.tex
%    pdflatex rapport-estore.tex   (deux passes pour les références)
%
%  Ou plus simple : importer ce fichier sur Overleaf (overleaf.com)
% =====================================================================

\documentclass[11pt,a4paper,oneside]{report}

% --- Encodage et langue ---
\usepackage[utf8]{inputenc}
\usepackage[T1]{fontenc}
\usepackage[french]{babel}
\usepackage{lmodern}
\usepackage{microtype}

% --- Mise en page ---
\usepackage[a4paper,top=2.5cm,bottom=2.5cm,left=2.5cm,right=2.5cm]{geometry}
\usepackage{setspace}
\onehalfspacing
\usepackage{parskip}
\setlength{\parindent}{0pt}
\setlength{\parskip}{0.6em}

% --- Couleurs ---
\usepackage{xcolor}
\definecolor{primaryblue}{HTML}{0D6EFD}
\definecolor{darknavy}{HTML}{0A2540}
\definecolor{accentgreen}{HTML}{198754}
\definecolor{accentorange}{HTML}{F77F00}
\definecolor{lightgray}{HTML}{F5F7FA}
\definecolor{borderlight}{HTML}{DEE2E6}
\definecolor{codebg}{HTML}{F8F9FA}
\definecolor{codecomment}{HTML}{6A737D}
\definecolor{codekeyword}{HTML}{D73A49}
\definecolor{codestring}{HTML}{032F62}
\definecolor{codenumber}{HTML}{005CC5}

% --- Graphismes et tableaux ---
\usepackage{graphicx}
\usepackage{booktabs}
\usepackage{longtable}
\usepackage{array}
\usepackage{multirow}
\usepackage{tabularx}
\newcolumntype{L}[1]{>{\raggedright\arraybackslash}p{#1}}
\usepackage{caption}
\captionsetup[figure]{labelfont={bf,color=primaryblue}, textfont=it, font=small}
\captionsetup[table]{labelfont={bf,color=primaryblue}, textfont=it, font=small}

% --- Listings (code) ---
\usepackage{listings}
\lstset{
  backgroundcolor=\color{codebg},
  basicstyle=\ttfamily\footnotesize,
  breakatwhitespace=false,
  breaklines=true,
  captionpos=b,
  commentstyle=\color{codecomment}\itshape,
  keywordstyle=\color{codekeyword}\bfseries,
  stringstyle=\color{codestring},
  numberstyle=\tiny\color{codecomment},
  numbers=left,
  numbersep=10pt,
  showspaces=false,
  showstringspaces=false,
  showtabs=false,
  tabsize=2,
  frame=leftline,
  framesep=8pt,
  rulecolor=\color{primaryblue},
  framerule=2pt,
  xleftmargin=20pt,
  xrightmargin=5pt,
  belowskip=1em,
  aboveskip=1em
}

\lstdefinelanguage{Java2}{
  morekeywords={var,record,sealed,permits,yield},
  keywordstyle=\color{codekeyword}\bfseries,
  language=Java
}

\lstdefinelanguage{TypeScript}{
  morekeywords={
    abstract,as,async,await,break,case,catch,class,const,continue,declare,
    default,delete,do,else,enum,export,extends,finally,for,from,function,
    get,if,implements,import,in,instanceof,interface,is,let,module,namespace,
    new,of,package,private,protected,public,readonly,return,set,static,
    super,switch,throw,try,type,typeof,var,void,while,with,yield,
    string,number,boolean,any,unknown,never,null,undefined,this,
    inject,signal,computed,effect
  },
  sensitive=true,
  morecomment=[l]{//},
  morecomment=[s]{/*}{*/},
  morestring=[b]',
  morestring=[b]",
  morestring=[b]`,
  keywordstyle=\color{codekeyword}\bfseries,
  stringstyle=\color{codestring},
  commentstyle=\color{codecomment}\itshape
}

\lstdefinestyle{javastyle}{language=Java2, basicstyle=\ttfamily\footnotesize}
\lstdefinestyle{tsstyle}{language=TypeScript, basicstyle=\ttfamily\footnotesize}
\lstdefinestyle{xmlstyle}{language=XML, basicstyle=\ttfamily\footnotesize}
\lstdefinestyle{shellstyle}{language=bash, basicstyle=\ttfamily\footnotesize, keywordstyle=\color{codekeyword}}
\lstdefinestyle{propertiesstyle}{
  basicstyle=\ttfamily\footnotesize,
  morekeywords={spring,server,jwt,logging},
  keywordstyle=\color{primaryblue}\bfseries,
  commentstyle=\color{codecomment}\itshape,
  morecomment=[l]{\#}
}

% --- Boîtes colorées ---
\usepackage[most]{tcolorbox}
\newtcolorbox{infobox}[1][]{
  colback=primaryblue!8, colframe=primaryblue,
  boxrule=0.4pt, arc=2pt, left=8pt, right=8pt, top=6pt, bottom=6pt,
  fonttitle=\bfseries, title=\faIcon{info-circle}~Note,#1
}
\newtcolorbox{tipbox}[1][]{
  colback=accentgreen!8, colframe=accentgreen,
  boxrule=0.4pt, arc=2pt, left=8pt, right=8pt, top=6pt, bottom=6pt,
  fonttitle=\bfseries, title=Astuce,#1
}
\newtcolorbox{warnbox}[1][]{
  colback=accentorange!10, colframe=accentorange,
  boxrule=0.4pt, arc=2pt, left=8pt, right=8pt, top=6pt, bottom=6pt,
  fonttitle=\bfseries, title=Attention,#1
}

% --- Titres de chapitres et sections ---
\usepackage{titlesec}
\titleformat{\chapter}[display]
  {\normalfont\huge\bfseries\color{primaryblue}}
  {\filleft\Large CHAPITRE \thechapter}
  {1ex}
  {\titlerule[1.5pt]\vspace{1ex}\filleft}
  [\vspace{0.5ex}\titlerule[0.5pt]]

\titleformat{\section}
  {\Large\bfseries\color{darknavy}}
  {\thesection}{1em}{}[{\titlerule[0.4pt]}]

\titleformat{\subsection}
  {\large\bfseries\color{primaryblue}}
  {\thesubsection}{1em}{}

\titleformat{\subsubsection}
  {\normalsize\bfseries\color{darknavy}}
  {\thesubsubsection}{1em}{}

% --- En-têtes et pieds de page ---
\usepackage{fancyhdr}
\pagestyle{fancy}
\fancyhf{}
\fancyhead[L]{\small\textcolor{primaryblue}{\textsc{E-Store}}}
\fancyhead[R]{\small\textit{\leftmark}}
\fancyfoot[L]{\small\textcolor{codecomment}{Belmoussa \& Ben Soumane — 2025-2026}}
\fancyfoot[R]{\small\textbf{\thepage}}
\renewcommand{\headrulewidth}{0.4pt}
\renewcommand{\footrulewidth}{0.2pt}

\fancypagestyle{plain}{
  \fancyhf{}
  \fancyfoot[L]{\small\textcolor{codecomment}{Belmoussa \& Ben Soumane — 2025-2026}}
  \fancyfoot[R]{\small\textbf{\thepage}}
  \renewcommand{\headrulewidth}{0pt}
  \renewcommand{\footrulewidth}{0.2pt}
}

% --- Hyperliens ---
\usepackage[colorlinks=true,linkcolor=primaryblue,urlcolor=primaryblue,
            citecolor=primaryblue,filecolor=primaryblue,
            pdftitle={Rapport E-Store},
            pdfauthor={Akram Belmoussa & Nouhaila Ben Soumane}]{hyperref}

% --- TikZ pour les schémas ---
\usepackage{tikz}
\usetikzlibrary{shapes,arrows,positioning,fit,backgrounds,decorations.pathreplacing,
                shapes.geometric,calc,arrows.meta}

% --- Liste des abréviations ---
\usepackage{enumitem}

% =====================================================================
%  DOCUMENT
% =====================================================================
\begin{document}

% --- PAGE DE GARDE ---
\begin{titlepage}
  \begin{center}
  \begin{tikzpicture}[remember picture, overlay]
    \fill[primaryblue] ([yshift=-2cm]current page.north west) rectangle ([yshift=-2.3cm]current page.north east);
    \fill[primaryblue] ([yshift=2.3cm]current page.south west) rectangle ([yshift=2cm]current page.south east);
  \end{tikzpicture}

  \vspace*{1cm}
  {\large\textsc{Royaume du Maroc}}\\[0.2cm]
  {\large\textsc{Université Hassan II de Casablanca}}\\[0.2cm]
  {\large\textsc{Faculté des Sciences Ben M'Sick}}\\[0.6cm]

  \rule{\linewidth}{1pt}\\[0.3cm]
  {\Large \textbf{Département d'Informatique}}\\[0.2cm]
  {\large Master/Licence — Module \textbf{Full-Stack}}\\
  \rule{\linewidth}{1pt}\\[2cm]

  \vfill
  {\Huge\bfseries\color{primaryblue} Rapport de Projet de Fin de Module}\\[1cm]
  \rule[0.5cm]{0.4\linewidth}{0.4pt}\\[0.3cm]
  {\Huge\bfseries\color{darknavy} E-STORE}\\[0.4cm]
  {\LARGE\itshape Conception et réalisation d'une plateforme}\\[0.2cm]
  {\LARGE\itshape e-commerce full-stack à persistance hybride}\\[0.5cm]
  \rule{0.4\linewidth}{0.4pt}\\[2cm]
  \vfill

  \begin{minipage}{0.46\textwidth}
    \begin{flushleft}
      \textbf{\large Réalisé par :}\\[0.2cm]
      Akram \textsc{Belmoussa}\\
      Nouhaila \textsc{Ben Soumane}
    \end{flushleft}
  \end{minipage}
  \hfill
  \begin{minipage}{0.46\textwidth}
    \begin{flushright}
      \textbf{\large Encadré par :}\\[0.2cm]
      Pr.\ Omar \textsc{Zahour}\\
      Faculté des Sciences Ben M'Sick
    \end{flushright}
  \end{minipage}

  \vspace{2cm}
  {\large Année universitaire \textbf{2025 -- 2026}}
  \end{center}
\end{titlepage}

% --- DÉDICACES ---
\thispagestyle{empty}
\vspace*{2cm}
\begin{flushright}\itshape\large
À nos parents, pour leur soutien indéfectible.\\[0.4cm]
À nos enseignants, pour le savoir partagé.\\[0.4cm]
À nos camarades de promotion, pour la collaboration.\\[0.4cm]
\end{flushright}
\vfill
\begin{center}\itshape
\textbf{Akram \& Nouhaila}
\end{center}
\newpage

% --- REMERCIEMENTS ---
\thispagestyle{plain}
\chapter*{Remerciements}
\addcontentsline{toc}{chapter}{Remerciements}

Au terme de ce projet, il nous est agréable d'exprimer nos sincères
remerciements à toutes les personnes qui ont contribué, de près ou de loin,
à sa réalisation.

Nous tenons en premier lieu à remercier notre encadrant pédagogique,
\textbf{Pr.\ Omar Zahour}, pour la qualité de son enseignement, sa
disponibilité, ses conseils avisés et la rigueur scientifique qu'il a su
nous transmettre tout au long de ce module. Ses orientations ont été
déterminantes dans la définition du périmètre du projet et dans les choix
techniques que nous avons opérés.

Nous remercions également l'ensemble du \textbf{corps professoral} de la
Faculté des Sciences Ben M'Sick — Université Hassan II de Casablanca, dont
les enseignements antérieurs en programmation orientée objet, en bases de
données, en génie logiciel et en architectures web ont constitué les
fondations sur lesquelles ce travail a pu s'élever.

Nos remerciements vont enfin à nos \textbf{familles} et à nos
\textbf{camarades de promotion}, pour leur soutien constant, leurs
relectures attentives et l'émulation intellectuelle qui nous a tirés vers
le haut.

% --- RÉSUMÉ ---
\chapter*{Résumé}
\addcontentsline{toc}{chapter}{Résumé}

Ce rapport présente la conception et la réalisation de \textbf{E-Store},
une application e-commerce complète développée dans le cadre du module
\textit{Full-Stack} de la Faculté des Sciences Ben M'Sick. Le projet adopte
une architecture rigoureuse en \textbf{trois couches techniques}
(présentation, logique métier, accès aux données) et un découpage en
\textbf{cinq domaines fonctionnels} suivant les principes du
\textit{Domain-Driven Design}.

L'application implémente l'ensemble du parcours utilisateur d'un site
marchand : inscription, authentification, navigation dans un catalogue,
gestion d'un panier, validation transactionnelle de commandes et dépôt
d'avis produits. La particularité technique majeure réside dans la mise en
\oe uvre d'une \textbf{persistance hybride} associant un système relationnel
(MySQL ou H2) pour les données transactionnelles à un système documentaire
(MongoDB) pour les avis utilisateurs.

Le backend repose sur \textbf{Spring Boot 3.3} et \textbf{Spring Security 6}
avec une authentification stateless par jeton \textbf{JWT}. Le frontend
exploite les nouveautés d'\textbf{Angular 17} (composants \textit{standalone},
\textit{signals}, lazy loading) et le framework \textit{Bootstrap 5}. Une
suite de neuf tests unitaires JUnit 5 / Mockito couvre les services métier
critiques.

\textbf{Mots-clés :} Spring Boot, Angular, JWT, MongoDB, Architecture en
couches, Domain-Driven Design, Persistance hybride, ACID, REST.

% --- TABLE DES MATIÈRES ---
\tableofcontents
\listoffigures
\listoftables

% --- ABRÉVIATIONS ---
\chapter*{Liste des abréviations}
\addcontentsline{toc}{chapter}{Liste des abréviations}

\begin{description}[leftmargin=4cm,style=nextline]
  \item[ACID] Atomicity, Consistency, Isolation, Durability
  \item[API] Application Programming Interface
  \item[BCrypt] Blowfish Crypt (algorithme de hachage)
  \item[CRUD] Create, Read, Update, Delete
  \item[CORS] Cross-Origin Resource Sharing
  \item[DDD] Domain-Driven Design
  \item[DTO] Data Transfer Object
  \item[ERD] Entity Relationship Diagram
  \item[FSBM] Faculté des Sciences Ben M'Sick
  \item[H2] Base de données relationnelle in-memory en Java
  \item[HTTP] HyperText Transfer Protocol
  \item[IDE] Integrated Development Environment
  \item[JPA] Jakarta Persistence API
  \item[JSON] JavaScript Object Notation
  \item[JVM] Java Virtual Machine
  \item[JWT] JSON Web Token
  \item[MCD] Modèle Conceptuel de Données
  \item[NoSQL] Not Only SQL
  \item[ORM] Object-Relational Mapping
  \item[REST] Representational State Transfer
  \item[SGBD] Système de Gestion de Bases de Données
  \item[SPA] Single Page Application
  \item[SQL] Structured Query Language
  \item[UI/UX] User Interface / User eXperience
  \item[XML] eXtensible Markup Language
\end{description}

% =====================================================================
%  INTRODUCTION GÉNÉRALE
% =====================================================================
\chapter*{Introduction générale}
\addcontentsline{toc}{chapter}{Introduction générale}
\markboth{INTRODUCTION GÉNÉRALE}{}

L'avènement du commerce électronique a profondément remodelé les pratiques
commerciales mondiales. Selon les estimations de la
\textit{Confédération Marocaine du Commerce Électronique}, le volume des
transactions en ligne au Maroc a franchi le cap des \textbf{20 milliards de
dirhams} en 2024, avec une croissance annuelle moyenne supérieure à 25\%.
Cette dynamique s'accompagne d'une exigence accrue en termes
d'\textit{ergonomie}, de \textit{performance} et de \textit{sécurité} des
plateformes marchandes.

Le module \textbf{Full-Stack}, dispensé sous la direction du
\textbf{Pr.\ Omar Zahour}, vise à doter les étudiants des compétences
nécessaires à la conception et à la réalisation de telles applications
\textit{de bout en bout} — du modèle de données jusqu'à l'interface
utilisateur, en passant par la logique métier, la sécurité et le
déploiement.

C'est dans ce contexte que s'inscrit le présent projet, baptisé
\textbf{E-Store}. Il s'agit d'une mise en pratique exhaustive des
enseignements théoriques par la construction d'une application e-commerce
complète, intégrant les meilleurs standards de l'industrie en 2025 :
Spring Boot 3.3, Angular 17, persistance hybride relationnel/documentaire,
authentification JWT, tests unitaires automatisés et conteneurisation
Docker.

\section*{Objectifs du projet}

Ce projet poursuit cinq objectifs principaux :

\begin{enumerate}
  \item \textbf{Maîtriser une architecture full-stack moderne} en
    appliquant un découpage en couches et en domaines fonctionnels rigoureux.
  \item \textbf{Illustrer la persistance hybride} en faisant cohabiter, au
    sein d'une même application, un SGBD relationnel (MySQL) et un SGBD
    documentaire (MongoDB).
  \item \textbf{Appliquer les bonnes pratiques} de génie logiciel : DTOs,
    transactions, validation, gestion centralisée des erreurs, tests unitaires.
  \item \textbf{Sécuriser l'application} par un mécanisme stateless reposant
    sur Spring Security 6, BCrypt et JWT.
  \item \textbf{Produire une interface utilisateur réactive} avec Angular 17
    en mode \textit{standalone}, Bootstrap 5 et signals.
\end{enumerate}

\section*{Plan du document}

Le présent rapport s'articule autour de \textbf{neuf chapitres} :

\begin{itemize}[leftmargin=*]
  \item Le \textbf{chapitre 1} pose le contexte du projet et formalise le
    cahier des charges fonctionnel et non fonctionnel.
  \item Le \textbf{chapitre 2} décrit l'architecture générale, le découpage
    en couches et en domaines.
  \item Le \textbf{chapitre 3} justifie les choix techniques opérés pour
    chaque brique de la stack.
  \item Les \textbf{chapitres 4, 5 et 6} détaillent respectivement la
    réalisation du backend, du module de sécurité et du frontend.
  \item Le \textbf{chapitre 7} expose la stratégie de tests et les résultats
    obtenus.
  \item Le \textbf{chapitre 8} traite du déploiement, de la conteneurisation
    et de l'intégration continue.
  \item Le \textbf{chapitre 9} relate les difficultés rencontrées et les
    solutions mises en \oe uvre.
  \item Une \textbf{conclusion générale} dresse le bilan du travail
    accompli et trace les perspectives d'évolution.
\end{itemize}

% =====================================================================
%  CHAPITRE 1 : CONTEXTE ET CAHIER DES CHARGES
% =====================================================================
\chapter{Contexte et cahier des charges}

\section{Contexte académique}

Ce projet est réalisé dans le cadre du module \textit{Full-Stack} du
département d'informatique de la Faculté des Sciences Ben M'Sick. Ce module
intervient en fin de cursus, après que les étudiants ont acquis les bases
en \textit{programmation orientée objet} (Java), \textit{bases de données
relationnelles} (SQL, normalisation, MERISE), \textit{développement web
côté client} (HTML, CSS, JavaScript) et \textit{génie logiciel}
(UML, patrons de conception).

L'objectif pédagogique est triple :
\begin{enumerate}
  \item \textbf{Intégrer} les acquis des modules précédents au sein d'un
    projet de taille significative.
  \item \textbf{Découvrir} les frameworks de production utilisés dans
    l'industrie (Spring Boot, Angular).
  \item \textbf{Pratiquer} le travail en binôme sur un dépôt Git partagé,
    en respectant des conventions de nommage et de commit.
\end{enumerate}

\section{Étude du marché et de l'existant}

\subsection{Marché du commerce électronique au Maroc}

Le commerce électronique marocain connaît une expansion soutenue.
Les plateformes leaders du marché — \textit{Jumia}, \textit{Avito},
\textit{Marjane.ma}, \textit{Decathlon.ma} — partagent un certain nombre de
\textbf{fonctionnalités-types} qui constituent l'état de l'art :
authentification utilisateur, navigation dans un catalogue catégorisé,
recherche textuelle et par filtres, gestion d'un panier persistant,
processus de commande sécurisé, suivi d'historique, dépôt et consultation
d'avis clients, gestion de favoris.

\subsection{Forces et faiblesses des plateformes existantes}

\begin{table}[htbp]
\centering
\caption{Analyse comparative succincte de plateformes existantes}
\label{tab:concurrents}
\small
\begin{tabularx}{\textwidth}{lXX}
\toprule
\textbf{Critère} & \textbf{Forces communes} & \textbf{Faiblesses récurrentes}\\
\midrule
Catalogue   & Richesse, pagination     & Filtres parfois lents, pas toujours pertinents \\
Panier      & Persistance, modification quantité & Synchronisation multi-appareils inégale \\
Sécurité    & HTTPS, JWT/OAuth         & Mots de passe faibles autorisés \\
Performance & CDN, cache               & Surcharge en période de soldes \\
Avis        & Notation par étoiles     & Modération absente ou tardive \\
\bottomrule
\end{tabularx}
\end{table}

E-Store ne prétend évidemment pas rivaliser avec ces géants, mais propose
une \textit{maquette pédagogique fonctionnelle} qui en reprend les
fondamentaux dans un périmètre maîtrisable.

\section{Cahier des charges fonctionnel}

\subsection{Acteurs}

L'application distingue deux profils d'utilisateurs :

\begin{itemize}
  \item \textbf{USER} (utilisateur standard) : peut s'inscrire, se connecter,
    consulter le catalogue, ajouter au panier, passer commande, consulter
    son historique et déposer des avis.
  \item \textbf{ADMIN} (administrateur) : dispose, en plus, des droits de
    création, modification et suppression sur les catégories, les produits
    et le stock.
\end{itemize}

\subsection{Cas d'usage principaux}

\begin{table}[htbp]
\centering
\caption{Cas d'usage fonctionnels couverts par E-Store}
\label{tab:usecases}
\small
\begin{tabularx}{\textwidth}{p{3.5cm}lX}
\toprule
\textbf{Cas d'usage} & \textbf{Acteur} & \textbf{Description}\\
\midrule
Inscription   & Visiteur & Création d'un compte avec email unique et mot de passe ($\geq$6 car.)\\
Connexion     & Utilisateur & Authentification email/mot de passe, retour d'un JWT 24h\\
Profil        & USER & Consultation et mise à jour du profil (téléphone, adresse)\\
Catalogue     & Tous & Liste paginée, recherche textuelle, filtre par catégorie\\
Fiche produit & Tous & Détail (prix, stock, image), liste des avis\\
Panier        & USER & Ajout, modification, suppression, vidage\\
Commande      & USER & Validation transactionnelle (ACID), historique\\
Avis          & USER & Dépôt (note 1-5 + commentaire), consultation\\
Catégorie/Produit & ADMIN & CRUD complet\\
Stock         & ADMIN & Lecture publique, mise à jour réservée\\
\bottomrule
\end{tabularx}
\end{table}

\section{Cahier des charges non fonctionnel}

\subsection{Performance}
Les pages doivent se charger en moins de \textbf{2 secondes} sur connexion
nominale (4G ou Wi-Fi). Le \textit{bundle initial} JavaScript doit rester
sous \textbf{50 ko} grâce au lazy loading.

\subsection{Sécurité}
\begin{itemize}
  \item Mots de passe hachés en \textbf{BCrypt} (jamais en clair).
  \item Authentification stateless par \textbf{JWT signé HMAC-SHA256},
    expiration \textbf{24 heures}.
  \item Endpoints d'administration protégés par \texttt{@PreAuthorize}.
  \item Validation systématique des entrées (\textit{Bean Validation}).
  \item CORS restreint à l'origine \texttt{localhost:4200} en développement.
\end{itemize}

\subsection{Maintenabilité}
\begin{itemize}
  \item Code conforme aux conventions Java/TypeScript.
  \item Découpage en cinq domaines fonctionnels.
  \item Au moins \textbf{neuf tests unitaires} couvrant les services métier
    critiques.
  \item Documentation interne (Javadoc, README, commentaires).
\end{itemize}

\subsection{Portabilité}
\begin{itemize}
  \item Démarrage possible \textbf{sans installation préalable} grâce au
    profil \texttt{dev} (H2 in-memory).
  \item Démarrage Docker en une commande pour le profil \texttt{prod}.
  \item Compatibilité multi-OS (Windows, macOS, Linux).
\end{itemize}

\section{Méthodologie de travail}

Nous avons retenu une approche \textbf{itérative légère}, inspirée d'Agile
mais adaptée au contexte d'un binôme étudiant. Le projet a été découpé en
\textbf{quinze étapes incrémentales}, chacune se concluant par un commit
Git significatif. La répartition des rôles a été la suivante :

\begin{itemize}
  \item \textbf{Akram Belmoussa} — Backend (Spring Boot, JPA, Sécurité JWT),
    intégration MongoDB, tests unitaires.
  \item \textbf{Nouhaila Ben Soumane} — Frontend (Angular 17, Bootstrap),
    modélisation des données, documentation et rapport.
  \item \textbf{Travail commun} — Architecture, design API, débogage,
    revue de code.
\end{itemize}

% =====================================================================
%  CHAPITRE 2 : ARCHITECTURE ET CONCEPTION
% =====================================================================
\chapter{Architecture et conception}

\section{Architecture générale en trois couches}

E-Store adopte une architecture \textbf{n-tiers classique} en trois couches,
qui présente l'avantage de séparer clairement les responsabilités et de
faciliter la maintenance. La figure~\ref{fig:archi-3layers} en donne une
vue synthétique.

\begin{figure}[htbp]
\centering
\begin{tikzpicture}[node distance=0.6cm, every node/.style={font=\small}]
  \tikzstyle{layer} = [rectangle, rounded corners=5pt, draw=primaryblue,
                       fill=primaryblue!10, text width=11cm, minimum height=1.5cm,
                       align=center, inner sep=8pt]
  \tikzstyle{tech} = [font=\itshape\small, color=darknavy]
  \tikzstyle{arrow} = [-{Stealth}, thick, color=primaryblue]

  \node[layer] (presentation) {%
    \textbf{Couche Présentation}\\
    \tech{Angular 17 standalone — Bootstrap 5 — TypeScript — RxJS / Signals}};
  \node[layer, below=of presentation] (business) {%
    \textbf{Couche Logique métier}\\
    \tech{Spring Boot 3.3 — Controllers REST — Services — DTOs — Spring Security 6}};
  \node[layer, below=of business] (data) {%
    \textbf{Couche Accès aux données}\\
    \tech{Spring Data JPA — Spring Data MongoDB — MySQL 8 / H2 — MongoDB 7}};

  \draw[arrow] (presentation) -- node[right, tech] {HTTP / JSON} (business);
  \draw[arrow] (business) -- node[right, tech] {SQL / BSON} (data);
\end{tikzpicture}
\caption{Architecture en trois couches de l'application E-Store}
\label{fig:archi-3layers}
\end{figure}

\subsection{Couche présentation (frontend)}
Cette couche est en charge de l'interaction avec l'utilisateur final. Elle
n'a aucune connaissance du modèle de persistance et communique avec le
backend exclusivement via une \textbf{API REST/JSON}.

\subsection{Couche logique métier (backend)}
C'est le c\oe ur applicatif. Elle reçoit les requêtes HTTP, valide les
entrées, applique les règles de gestion (vérification de stock,
calcul de totaux, génération de jetons), orchestre les transactions et
renvoie des DTOs sérialisés en JSON.

\subsection{Couche d'accès aux données}
Cette couche, masquée derrière les abstractions Spring Data, gère la
persistance effective. Deux SGBD coexistent :
\begin{itemize}
  \item \textbf{MySQL 8} (ou H2 en mode dev) pour les données strictement
    relationnelles : utilisateurs, produits, panier, commandes.
  \item \textbf{MongoDB 7} pour les avis, données semi-structurées en
    forte croissance.
\end{itemize}

\section{Découpage en domaines fonctionnels}

Chaque couche est elle-même découpée en \textbf{cinq domaines fonctionnels}
cohésifs, suivant les principes du \textit{Domain-Driven Design}
(figure~\ref{fig:domains}). Un sixième domaine, \texttt{review}, gère la
particularité MongoDB.

\begin{figure}[htbp]
\centering
\begin{tikzpicture}[node distance=0.4cm and 0.6cm, every node/.style={font=\small}]
  \tikzstyle{domain} = [rectangle, rounded corners=4pt, draw=primaryblue,
                        fill=lightgray, text width=4cm, minimum height=2.2cm,
                        align=center, inner sep=4pt]
  \tikzstyle{title} = [font=\bfseries\small, color=primaryblue]

  \node[domain] (customer) {%
    \title{customer}\\\smallskip
    User, Profile,\\Auth, JWT};
  \node[domain, right=of customer] (catalog) {%
    \title{catalog}\\\smallskip
    Category,\\Product, recherche};
  \node[domain, right=of catalog] (inventory) {%
    \title{inventory}\\\smallskip
    Stock par\\produit};

  \node[domain, below=of customer] (shopping) {%
    \title{shopping}\\\smallskip
    Cart,\\CartItem};
  \node[domain, below=of catalog] (billing) {%
    \title{billing}\\\smallskip
    Order, OrderItem\\@Transactional};
  \node[domain, below=of inventory] (review) {%
    \title{review}\\\smallskip
    Avis (MongoDB)};
\end{tikzpicture}
\caption{Les six domaines fonctionnels de l'application}
\label{fig:domains}
\end{figure}

Cette double organisation \textit{couches × domaines} produit la matrice
suivante (table~\ref{tab:matrix}) qui donne une vue exhaustive du projet.

\begin{table}[htbp]
\centering
\caption{Matrice architecturale : 3 couches × 6 domaines}
\label{tab:matrix}
\small
\begin{tabularx}{\textwidth}{lXXX}
\toprule
\textbf{Domaine} & \textbf{Présentation} & \textbf{Logique métier} & \textbf{Données}\\
\midrule
customer  & \texttt{login.component.ts}  & \texttt{AuthService}, JwtFilter   & \texttt{users}, \texttt{profiles} (JPA)\\
catalog   & \texttt{product-list}        & \texttt{ProductService}, recherche & \texttt{products}, \texttt{categories}\\
inventory & Indicateur stock UI         & \texttt{InventoryService}         & \texttt{inventories}\\
shopping  & \texttt{cart.component}      & \texttt{CartService}              & \texttt{carts}, \texttt{cart\_items}\\
billing   & \texttt{orders.component}    & \texttt{OrderService} \texttt{@Transactional} & \texttt{orders}, \texttt{order\_items}\\
review    & \texttt{review-form}         & \texttt{ReviewService}            & MongoDB \texttt{reviews}\\
\bottomrule
\end{tabularx}
\end{table}

\section{Modèle conceptuel des données relationnel}

Le modèle relationnel comporte \textbf{neuf entités} reliées par des
associations \textit{one-to-one}, \textit{one-to-many} et
\textit{many-to-one}. La figure~\ref{fig:erd} en présente la version
simplifiée.

\begin{figure}[htbp]
\centering
\begin{tikzpicture}[node distance=0.4cm and 0.5cm, every node/.style={font=\footnotesize}]
  \tikzstyle{entity} = [rectangle split, rectangle split parts=2, draw=primaryblue,
                        fill=lightgray, rounded corners=2pt, inner sep=3pt,
                        rectangle split part fill={primaryblue!20, lightgray}]
  \tikzstyle{rel} = [-{Stealth}, thick, color=darknavy]

  \node[entity] (user) {\textbf{User} \nodepart{second}
    id, firstName, lastName,\\email, password, role};
  \node[entity, below right=of user] (profile) {\textbf{Profile}\nodepart{second}
    id, phone, address,\\city, country};
  \node[entity, above right=of user] (cart) {\textbf{Cart}\nodepart{second}
    id, createdAt, updatedAt};
  \node[entity, right=4cm of cart] (cartitem) {\textbf{CartItem}\nodepart{second}
    id, quantity, unitPrice};

  \node[entity, below=2cm of cartitem] (product) {\textbf{Product}\nodepart{second}
    id, name, description,\\price, imageUrl};
  \node[entity, below right=of product] (inventory) {\textbf{Inventory}\nodepart{second}
    id, quantity};
  \node[entity, above right=of product] (category) {\textbf{Category}\nodepart{second}
    id, name, description};

  \node[entity, below=of user] (order) {\textbf{Order}\nodepart{second}
    id, orderDate, totalAmount, status};
  \node[entity, right=of order] (orderitem) {\textbf{OrderItem}\nodepart{second}
    id, quantity, unitPrice};

  \draw[rel] (user) -- node[above, sloped, font=\tiny] {1..1} (profile);
  \draw[rel] (user) -- node[above, sloped, font=\tiny] {1..1} (cart);
  \draw[rel] (user) -- node[above, sloped, font=\tiny] {1..n} (order);
  \draw[rel] (cart) -- node[above, font=\tiny] {1..n} (cartitem);
  \draw[rel] (cartitem) -- node[above, sloped, font=\tiny] {n..1} (product);
  \draw[rel] (product) -- node[above, font=\tiny] {1..1} (inventory);
  \draw[rel] (product) -- node[above, sloped, font=\tiny] {n..1} (category);
  \draw[rel] (order) -- node[above, font=\tiny] {1..n} (orderitem);
  \draw[rel] (orderitem) -- node[above, sloped, font=\tiny] {n..1} (product);
\end{tikzpicture}
\caption{Modèle Conceptuel de Données simplifié (relationnel)}
\label{fig:erd}
\end{figure}

\section{Modèle documentaire MongoDB}

La collection \texttt{reviews} stocke les avis utilisateurs sous forme de
documents JSON. La figure~\ref{fig:reviewdoc} montre la structure type.

\begin{figure}[htbp]
\centering
\begin{tcolorbox}[colback=codebg, colframe=primaryblue!50, sharp corners,
                  width=0.8\textwidth, halign=left]
\begin{verbatim}
{
  "_id":         ObjectId("..."),
  "productId":   12,
  "userId":      5,
  "authorName":  "Akram Belmoussa",
  "rating":      5,
  "comment":     "Excellent produit !",
  "createdAt":   ISODate("2025-04-22T14:30:00Z")
}
\end{verbatim}
\end{tcolorbox}
\caption{Structure d'un document de la collection \texttt{reviews}}
\label{fig:reviewdoc}
\end{figure}

Deux index secondaires sont créés via les annotations
\texttt{@Indexed} sur les champs \texttt{productId} et \texttt{userId},
afin d'accélérer les requêtes \texttt{findByProductId} et
\texttt{findByUserId}.

\section{Architecture API REST}

Toute communication frontend/backend transite par une API REST décrite
dans la table~\ref{tab:endpoints}. L'API respecte les standards :

\begin{itemize}
  \item Méthodes HTTP appropriées (GET / POST / PUT / DELETE).
  \item Codes de retour normalisés : 200, 201, 204, 400, 401, 403, 404, 409.
  \item Corps JSON encadré par une enveloppe \texttt{ApiResponse}
    standardisée.
  \item Versioning implicite par le préfixe \texttt{/api/}.
\end{itemize}

\begin{table}[htbp]
\centering
\caption{Endpoints REST principaux}
\label{tab:endpoints}
\small
\begin{tabularx}{\textwidth}{p{1.2cm}p{4cm}Xc}
\toprule
\textbf{Verbe} & \textbf{URL} & \textbf{Description} & \textbf{Auth.}\\
\midrule
POST   & /api/auth/register     & Inscription                & Public\\
POST   & /api/auth/login        & Connexion (renvoie JWT)    & Public\\
GET    & /api/users/me          & Profil utilisateur courant & USER\\
PUT    & /api/users/me          & Mise à jour du profil      & USER\\
GET    & /api/categories        & Liste des catégories       & Public\\
POST   & /api/categories        & Création (ADMIN)           & ADMIN\\
GET    & /api/products          & Catalogue + filtres        & Public\\
GET    & /api/products/\{id\}   & Fiche produit              & Public\\
POST   & /api/products          & Création                   & ADMIN\\
GET    & /api/inventory/\{id\}  & Lecture stock              & Public\\
PUT    & /api/inventory/\{id\}  & Modification stock         & ADMIN\\
GET    & /api/cart              & Panier courant             & USER\\
POST   & /api/cart/add          & Ajout au panier            & USER\\
DELETE & /api/cart/clear        & Vidage panier              & USER\\
POST   & /api/orders            & Validation commande        & USER\\
GET    & /api/orders            & Historique commandes       & USER\\
POST   & /api/reviews           & Dépôt d'avis               & USER\\
GET    & /api/reviews/product/\{id\} & Avis d'un produit     & Public\\
\bottomrule
\end{tabularx}
\end{table}

% =====================================================================
%  CHAPITRE 3 : CHOIX TECHNIQUES ET JUSTIFICATIONS
% =====================================================================
\chapter{Choix techniques et justifications}

\section{Backend : Java et Spring Boot}

\subsection{Pourquoi Java ?}
Java demeure, en 2025, l'un des langages les plus largement déployés en
entreprise — notamment dans le secteur bancaire, les télécommunications et
le e-commerce de grande échelle. Ses atouts pédagogiques sont multiples :
typage fort, programmation orientée objet stricte, écosystème mature,
machine virtuelle performante. La version \textbf{Java 17 LTS} (Long-Term
Support) a été retenue pour sa stabilité et son adoption massive.

\subsection{Pourquoi Spring Boot ?}
Spring Boot est le framework Java de référence pour la construction
d'applications backend modernes. Il automatise la configuration et fournit
une myriade de \textit{starters} prêts à l'emploi. Ses points forts sont :

\begin{itemize}
  \item \textbf{Auto-configuration} : par convention plutôt que par
    configuration explicite.
  \item \textbf{Embedded Tomcat} : déploiement en JAR auto-exécutable.
  \item \textbf{Spring Data} : abstraction puissante des SGBD.
  \item \textbf{Spring Security 6} : sécurisation standardisée.
  \item \textbf{Actuator} : observabilité prête à l'emploi.
\end{itemize}

La \textbf{version 3.3} a été choisie car compatible Jakarta EE 10 et
Java 17+, et stable depuis fin 2024.

\section{Persistance hybride : MySQL + MongoDB}

\subsection{Justification du choix hybride}
Le choix d'un \textbf{double système de persistance} n'est pas un caprice :
il répond à une réalité du terrain. Toutes les données ne sont pas égales :

\begin{itemize}
  \item Les \textbf{transactions financières} (commandes, panier, stock)
    exigent les propriétés ACID : atomicité, cohérence, isolation,
    durabilité. Le SQL relationnel y excelle.
  \item Les \textbf{avis utilisateurs} ne nécessitent ni jointures
    complexes, ni transactions ; ils croissent rapidement et sont
    consultés massivement en lecture. Le NoSQL documentaire les sert plus
    efficacement.
\end{itemize}

\begin{table}[htbp]
\centering
\caption{Comparaison MySQL vs MongoDB dans le contexte E-Store}
\label{tab:dbcompare}
\small
\begin{tabularx}{\textwidth}{lXX}
\toprule
\textbf{Critère} & \textbf{MySQL 8 (relationnel)} & \textbf{MongoDB 7 (documentaire)}\\
\midrule
Paradigme         & SQL, tables normalisées       & BSON, documents JSON\\
Schéma            & Strict, prédéfini             & Flexible\\
Transactions      & ACID natives                  & Limitées (multi-document v4+)\\
Joins             & Performants                   & Rares (\$lookup)\\
Lectures massives & Bonnes (index)                & Excellentes (réplicas)\\
Cas d'usage E-Store & Users, Products, Orders     & Reviews\\
\bottomrule
\end{tabularx}
\end{table}

\subsection{Bascule H2 / MySQL via les profils}
Le projet exploite la fonctionnalité Spring de \textbf{profils}
(\texttt{application-dev.properties}, \texttt{application-prod.properties})
pour basculer transparemment entre une base H2 in-memory (développement,
démonstration sans installation) et une base MySQL persistante
(production, soutenance avec données pérennes).

\begin{lstlisting}[style=propertiesstyle, caption=Extrait de application-dev.properties]
# Profil dev : H2 in-memory
spring.datasource.url=jdbc:h2:mem:estore;MODE=MySQL
spring.datasource.driver-class-name=org.h2.Driver
spring.jpa.database-platform=org.hibernate.dialect.H2Dialect
spring.h2.console.enabled=true
\end{lstlisting}

\section{Frontend : Angular 17 standalone}

\subsection{Pourquoi Angular ?}
Trois frameworks dominent le développement web SPA en 2025 :
\textbf{Angular}, \textbf{React} et \textbf{Vue}. Angular se distingue par
son approche \textit{tout-en-un} (router, formulaires, HTTP, i18n inclus) et
son langage TypeScript natif, qui apporte la rigueur du typage statique au
JavaScript. Pour un projet pédagogique de fin de module, Angular offre
l'environnement le plus structurant.

\subsection{Le mode standalone d'Angular 17}
Depuis la version 14, et de manière généralisée depuis la version 17,
Angular permet de s'affranchir des \texttt{NgModule} historiques au profit
de \textbf{composants standalone}. Avantages :

\begin{itemize}
  \item Code plus concis (suppression du boilerplate des modules).
  \item Lazy loading natif par route.
  \item Imports explicites par composant : meilleure traçabilité.
  \item Bundle initial plus léger.
\end{itemize}

\subsection{Signals : la nouvelle réactivité}
Les \textbf{signals} sont la grande nouveauté d'Angular 17. Ils
remplacent avantageusement, dans les cas simples, les \texttt{Observable}
RxJS pour la gestion d'état. Dans E-Store, ils servent par exemple à
afficher dynamiquement le badge du panier dans la barre de navigation :

\begin{lstlisting}[style=tsstyle, caption=Utilisation de signals dans CartService]
@Injectable({ providedIn: 'root' })
export class CartService {
  readonly cart = signal<Cart | null>(null);
  readonly itemCount = signal<number>(0);

  add(productId: number, qty: number) {
    return this.http.post(...).pipe(
      tap(r => this.update(r.data))
    );
  }
}
\end{lstlisting}

% =====================================================================
%  CHAPITRE 4 : RÉALISATION BACKEND
% =====================================================================
\chapter{Réalisation du backend}

\section{Structure générale du projet}

L'arborescence \texttt{estore-backend/src/main/java/com/estore/} reflète
fidèlement le découpage en domaines :

\begin{lstlisting}[style=shellstyle]
com.estore/
  EstoreApplication.java        <- Point d'entree
  customer/    {entity, dto, repository, service, controller, security}
  catalog/     {entity, dto, repository, service, controller}
  inventory/   {entity, dto, repository, service, controller}
  shopping/    {entity, dto, repository, service, controller}
  billing/     {entity, dto, repository, service, controller}
  review/      {document, dto, repository, service, controller}
  shared/      ApiResponse.java
  config/      CorsConfig, SecurityConfig, DataSeeder
  exception/   GlobalExceptionHandler + 3 exceptions metier
\end{lstlisting}

Chaque domaine est \textbf{autonome} : il ne peut dépendre que des couches
inférieures de son propre domaine, ou des composants partagés
(\texttt{shared}, \texttt{exception}). Les domaines communiquent
exclusivement via leurs services publics.

\section{Domaine \texttt{customer}}

Ce domaine pivote autour de l'entité \texttt{User}, à laquelle est associé
un \texttt{Profile} en relation \texttt{@OneToOne}.

\subsection{Entité User}

\begin{lstlisting}[style=javastyle, caption=Entite User (extrait)]
@Entity
@Table(name = "users")
public class User {
    @Id @GeneratedValue(strategy = GenerationType.IDENTITY)
    private Long id;

    @Column(nullable = false, unique = true)
    private String email;

    @JsonIgnore
    @Column(nullable = false)
    private String password;          // hash BCrypt, jamais en clair

    @Enumerated(EnumType.STRING)
    private Role role;                 // USER ou ADMIN

    @OneToOne(mappedBy = "user", cascade = CascadeType.ALL,
              orphanRemoval = true, fetch = FetchType.LAZY)
    private Profile profile;
}
\end{lstlisting}

L'annotation \texttt{@JsonIgnore} sur \texttt{password} empêche toute
sérialisation accidentelle. La cascade \texttt{ALL} sur \texttt{profile}
garantit l'atomicité création/suppression.

\subsection{AuthService}

L'\texttt{AuthService} centralise la logique d'inscription et de connexion.
La connexion délègue à Spring Security via l'\texttt{AuthenticationManager},
puis génère un JWT.

\begin{lstlisting}[style=javastyle, caption=Methode register de AuthService]
@Transactional
public AuthResponseDto register(RegisterDto dto) {
    if (userRepository.existsByEmail(dto.getEmail()))
        throw new BusinessException("Cet email est deja utilise");

    User user = User.builder()
        .firstName(dto.getFirstName()).lastName(dto.getLastName())
        .email(dto.getEmail())
        .password(passwordEncoder.encode(dto.getPassword()))  // BCrypt
        .role(Role.USER).build();

    user.setProfile(Profile.builder().user(user).build());
    User saved = userRepository.save(user);

    String token = jwtService.generateToken(saved);
    return AuthResponseDto.builder()
        .token(token).user(UserDto.from(saved)).build();
}
\end{lstlisting}

\section{Domaine \texttt{catalog}}

Le catalogue regroupe \texttt{Category} et \texttt{Product}. Le repository
\texttt{ProductRepository} expose une méthode personnalisée combinant
recherche textuelle et filtre par catégorie :

\begin{lstlisting}[style=javastyle, caption=Recherche multi-critere de produits]
public interface ProductRepository extends JpaRepository<Product, Long> {
    @Query("""
        SELECT p FROM Product p
        WHERE (:categoryId IS NULL OR p.category.id = :categoryId)
          AND (:q IS NULL OR LOWER(p.name) LIKE LOWER(CONCAT('%',:q,'%'))
                          OR LOWER(p.description) LIKE LOWER(CONCAT('%',:q,'%')))
    """)
    Page<Product> search(@Param("categoryId") Long categoryId,
                         @Param("q") String q, Pageable pageable);
}
\end{lstlisting}

L'utilisation de JPQL plutôt que de SQL natif garantit la portabilité
H2/MySQL, et la pagination via \texttt{Pageable} évite le retour massif
d'enregistrements.

\section{Domaine \texttt{shopping} et \texttt{billing}}

\subsection{Le panier (shopping)}
L'entité \texttt{Cart} maintient une relation \texttt{@OneToMany} avec
\texttt{CartItem}. À l'ajout d'un produit déjà présent, la quantité est
incrémentée plutôt que dupliquée.

\subsection{Le checkout transactionnel (billing)}
La validation de commande constitue \textbf{l'opération la plus critique}
de l'application. Elle touche trois agrégats (Cart, Order, Inventory) et
doit être strictement atomique.

\begin{lstlisting}[style=javastyle, caption=OrderService.checkout() — l'operation transactionnelle]
@Transactional
public OrderDto checkout() {
    User user = securityUtils.currentUser();
    Cart cart = cartService.getOrCreateCart(user);
    if (cart.getItems().isEmpty())
        throw new BusinessException("Votre panier est vide");

    // 1) Verifier le stock pour TOUS les items AVANT toute modification
    for (CartItem ci : cart.getItems())
        inventoryService.checkAvailability(
            ci.getProduct().getId(), ci.getQuantity());

    // 2) Creer la commande + decrementer le stock
    Order order = Order.builder().user(user)
        .status(OrderStatus.PENDING).build();
    BigDecimal total = BigDecimal.ZERO;
    for (CartItem ci : cart.getItems()) {
        order.getItems().add(OrderItem.builder()
            .order(order).product(ci.getProduct())
            .quantity(ci.getQuantity()).unitPrice(ci.getUnitPrice())
            .build());
        total = total.add(ci.getUnitPrice()
                            .multiply(BigDecimal.valueOf(ci.getQuantity())));
        inventoryService.decrement(
            ci.getProduct().getId(), ci.getQuantity());
    }
    order.setTotalAmount(total);
    order.setStatus(OrderStatus.CONFIRMED);

    // 3) Sauvegarder + vider le panier
    Order saved = orderRepository.save(order);
    cartService.clearCart(cart);
    return OrderDto.from(saved);
}
\end{lstlisting}

\begin{infobox}
L'annotation \texttt{@Transactional} de Spring garantit qu'en cas
d'exception levée à n'importe quelle étape (par exemple une rupture de
stock détectée tardivement, une coupure de la base de données, etc.), \emph{toutes} les modifications faites depuis l'entrée dans la méthode
sont automatiquement annulées (\textit{rollback}). C'est l'\textbf{atomicité
ACID} en action.
\end{infobox}

\section{Domaine \texttt{review} (MongoDB)}

Le document \texttt{Review} est mappé sur la collection MongoDB éponyme :

\begin{lstlisting}[style=javastyle, caption=Document MongoDB Review]
@Document(collection = "reviews")
public class Review {
    @Id private String id;
    @Indexed private Long productId;
    @Indexed private Long userId;
    private String authorName;
    private int rating;
    private String comment;
    private Instant createdAt;
}
\end{lstlisting}

Le repository tire profit de la \textit{convention over configuration} de
Spring Data : la simple déclaration d'une méthode au nom évocateur produit
la requête correspondante.

\begin{lstlisting}[style=javastyle]
public interface ReviewRepository extends MongoRepository<Review, String> {
    List<Review> findByProductIdOrderByCreatedAtDesc(Long productId);
    List<Review> findByUserIdOrderByCreatedAtDesc(Long userId);
}
\end{lstlisting}

\section{Gestion globale des erreurs}

Plutôt que de laisser remonter des exceptions Java brutes vers le client,
l'application centralise leur traitement via un
\texttt{@RestControllerAdvice}. Chaque exception métier est convertie en
réponse JSON propre avec le code HTTP approprié.

\begin{lstlisting}[style=javastyle, caption=Extrait de GlobalExceptionHandler]
@RestControllerAdvice
public class GlobalExceptionHandler {

    @ExceptionHandler(ResourceNotFoundException.class)
    public ResponseEntity<ApiResponse<Void>> handleNotFound(
            ResourceNotFoundException ex) {
        return ResponseEntity.status(HttpStatus.NOT_FOUND)
            .body(ApiResponse.error(ex.getMessage()));
    }

    @ExceptionHandler(BusinessException.class)
    public ResponseEntity<ApiResponse<Void>> handleBusiness(
            BusinessException ex) {
        return ResponseEntity.status(HttpStatus.CONFLICT)
            .body(ApiResponse.error(ex.getMessage()));
    }
    // ... + ValidationException, BadCredentialsException, AccessDeniedException
}
\end{lstlisting}

% =====================================================================
%  CHAPITRE 5 : SÉCURITÉ
% =====================================================================
\chapter{Sécurité de l'application}

La sécurité de toute application e-commerce est une exigence absolue.
E-Store met en \oe uvre une chaîne complète depuis le hachage des mots de
passe jusqu'à l'autorisation fine par rôle.

\section{Hachage BCrypt des mots de passe}

Aucun mot de passe utilisateur n'est jamais stocké en clair. Tous transitent
par l'algorithme \textbf{BCrypt}, un \textit{one-way hash} avec sel
aléatoire et facteur de coût ajustable. Le facteur 10 (par défaut Spring)
est utilisé, soit $2^{10} = 1024$ rounds.

\begin{lstlisting}[style=javastyle]
@Bean
public PasswordEncoder passwordEncoder() {
    return new BCryptPasswordEncoder();
}
\end{lstlisting}

\section{Authentification stateless par JWT}

L'authentification suit un modèle \textit{stateless} : aucun état n'est
conservé côté serveur entre deux requêtes. Le client porte lui-même la
preuve de son identité, sous la forme d'un \textbf{JSON Web Token} signé.

\subsection{Structure d'un JWT}
Un JWT est composé de trois parties séparées par un point :
\texttt{header.payload.signature}. Le payload de E-Store contient les
\textit{claims} suivants :

\begin{itemize}
  \item \texttt{sub} : email de l'utilisateur (subject)
  \item \texttt{uid} : identifiant numérique
  \item \texttt{role} : USER ou ADMIN
  \item \texttt{name} : nom complet (pour affichage UI)
  \item \texttt{iat} et \texttt{exp} : timestamps de création et d'expiration
\end{itemize}

\subsection{Cycle de vie complet}

La figure~\ref{fig:jwtflow} décrit le cycle de vie complet d'un JWT, de
sa génération à sa validation par le filtre Spring Security.

\begin{figure}[htbp]
\centering
\begin{tikzpicture}[node distance=0.6cm and 1.5cm, every node/.style={font=\footnotesize}]
  \tikzstyle{actor} = [rectangle, rounded corners=3pt, draw=primaryblue,
                       fill=primaryblue!15, text width=2.7cm, minimum height=0.8cm,
                       align=center]
  \tikzstyle{step} = [rectangle, rounded corners=3pt, draw=darknavy,
                      fill=lightgray, text width=8cm, minimum height=0.8cm,
                      align=left, inner sep=4pt]
  \tikzstyle{arrow} = [-{Stealth}, thick]

  \node[actor] (client) {Client (Angular)};
  \node[actor, right=4cm of client] (server) {Serveur (Spring)};

  \node[step, below=of client, xshift=2cm] (s1)
    {\textbf{1.} Login : POST /api/auth/login (email, password)};
  \node[step, below=of s1] (s2)
    {\textbf{2.} AuthService valide via passwordEncoder.matches()};
  \node[step, below=of s2] (s3)
    {\textbf{3.} JwtService.generateToken() — HMAC-SHA256, exp 24h};
  \node[step, below=of s3] (s4)
    {\textbf{4.} Client persiste le token dans localStorage};
  \node[step, below=of s4] (s5)
    {\textbf{5.} Toute requete : header Authorization: Bearer <token>};
  \node[step, below=of s5] (s6)
    {\textbf{6.} JwtAuthenticationFilter valide la signature et l'exp.};
  \node[step, below=of s6] (s7)
    {\textbf{7.} Si OK : SecurityContext.setAuthentication(...)};
\end{tikzpicture}
\caption{Cycle de vie d'un JWT dans E-Store}
\label{fig:jwtflow}
\end{figure}

\subsection{Génération du token}

\begin{lstlisting}[style=javastyle, caption=JwtService.generateToken]
public String generateToken(User user) {
    Map<String, Object> claims = new HashMap<>();
    claims.put("uid",  user.getId());
    claims.put("role", user.getRole().name());
    claims.put("name", user.getFirstName() + " " + user.getLastName());

    Date now = new Date();
    Date exp = new Date(now.getTime() + expirationMs);  // 24h

    return Jwts.builder()
        .claims(claims)
        .subject(user.getEmail())
        .issuedAt(now).expiration(exp)
        .signWith(key)                                  // HMAC-SHA256
        .compact();
}
\end{lstlisting}

\subsection{Validation côté serveur}

À chaque requête entrante, le \texttt{JwtAuthenticationFilter}
(extension de \texttt{OncePerRequestFilter}) extrait le token, vérifie sa
signature et son expiration, puis place l'utilisateur dans le
\texttt{SecurityContextHolder}.

\section{Autorisation par rôle}

Les opérations sensibles (création/modification/suppression de catégories
et de produits, mise à jour du stock) sont protégées par
\texttt{@PreAuthorize}.

\begin{lstlisting}[style=javastyle]
@PostMapping
@PreAuthorize("hasRole('ADMIN')")
public ResponseEntity<ApiResponse<ProductDto>> create(
        @Valid @RequestBody CreateProductDto dto) {
    return ResponseEntity.status(HttpStatus.CREATED)
        .body(ApiResponse.ok("Produit cree", productService.create(dto)));
}
\end{lstlisting}

Tout appel à cet endpoint par un utilisateur non-ADMIN se solde par un
HTTP 403 Forbidden, intercepté par le \texttt{GlobalExceptionHandler}.

\section{Configuration CORS}

En développement, le frontend tourne sur le port 4200 et le backend sur le
port 8080. Cette différence d'origine déclencherait une erreur CORS du
navigateur sans configuration appropriée.

\begin{lstlisting}[style=javastyle, caption=CorsConfig]
@Override
public void addCorsMappings(CorsRegistry registry) {
    registry.addMapping("/api/**")
        .allowedOrigins("http://localhost:4200")
        .allowedMethods("GET","POST","PUT","DELETE","OPTIONS")
        .allowedHeaders("*").allowCredentials(true);
}
\end{lstlisting}

% =====================================================================
%  CHAPITRE 6 : RÉALISATION FRONTEND
% =====================================================================
\chapter{Réalisation du frontend}

\section{Structure de l'application Angular}

Le projet Angular suit l'organisation \textbf{core / shared / features},
classique des applications Angular modernes :

\begin{lstlisting}[style=shellstyle]
src/app/
  app.component.ts          <- Composant racine
  app.config.ts             <- Configuration globale (providers)
  app.routes.ts             <- Routes
  core/
    services/               <- AuthService, CartService, ...
    guards/                 <- AuthGuard
    interceptors/           <- AuthInterceptor, ErrorInterceptor
    models/                 <- Interfaces TypeScript
  shared/
    components/             <- Header, Footer, Loader, Toast
  features/
    auth/                   <- Login, Register
    catalog/                <- ProductList, ProductDetail
    cart/                   <- Cart
    orders/                 <- Orders
    profile/                <- Profile
    reviews/                <- ReviewForm
\end{lstlisting}

\section{Bootstrap de l'application}

Le fichier \texttt{main.ts} amorce l'application sans \texttt{NgModule},
en passant directement par \texttt{bootstrapApplication}.

\begin{lstlisting}[style=tsstyle, caption=main.ts]
import { bootstrapApplication } from '@angular/platform-browser';
import { AppComponent } from './app/app.component';
import { appConfig } from './app/app.config';

bootstrapApplication(AppComponent, appConfig)
  .catch(err => console.error(err));
\end{lstlisting}

La configuration des providers globaux (router, HTTP client, intercepteurs)
est centralisée dans \texttt{app.config.ts} :

\begin{lstlisting}[style=tsstyle, caption=app.config.ts]
export const appConfig: ApplicationConfig = {
  providers: [
    provideZoneChangeDetection({ eventCoalescing: true }),
    provideRouter(routes, withComponentInputBinding()),
    provideHttpClient(withInterceptors([
      authInterceptor, errorInterceptor
    ]))
  ]
};
\end{lstlisting}

\section{Routing et lazy loading}

Toutes les routes utilisent \texttt{loadComponent} pour un lazy loading
natif, ce qui produit des bundles séparés par fonctionnalité.

\begin{lstlisting}[style=tsstyle, caption=app.routes.ts]
export const routes: Routes = [
  { path: '', loadComponent: () =>
      import('./features/catalog/product-list.component')
        .then(m => m.ProductListComponent) },
  { path: 'login', loadComponent: () =>
      import('./features/auth/login.component')
        .then(m => m.LoginComponent) },
  { path: 'cart', loadComponent: () =>
      import('./features/cart/cart.component')
        .then(m => m.CartComponent),
    canActivate: [authGuard] },
  // ...
];
\end{lstlisting}

Mesure : le bundle initial pèse \textbf{13.84 ko}, et chaque feature ajoute
en moyenne \textbf{8 ko} chargés à la demande.

\section{AuthInterceptor et ErrorInterceptor}

L'\texttt{AuthInterceptor} attache silencieusement le JWT à chaque requête
HTTP sortante :

\begin{lstlisting}[style=tsstyle]
export const authInterceptor: HttpInterceptorFn = (req, next) => {
  const token = inject(AuthService).token;
  if (!token) return next(req);
  return next(req.clone({
    setHeaders: { Authorization: `Bearer ${token}` }
  }));
};
\end{lstlisting}

L'\texttt{ErrorInterceptor} centralise la gestion d'erreurs :
401 entraîne la déconnexion automatique, 4xx/5xx affichent un toast.

\begin{lstlisting}[style=tsstyle]
export const errorInterceptor: HttpInterceptorFn = (req, next) => {
  const toast = inject(ToastService);
  const auth = inject(AuthService);
  const router = inject(Router);

  return next(req).pipe(
    catchError((err: HttpErrorResponse) => {
      if (err.status === 401) {
        auth.logout();
        router.navigate(['/login']);
        toast.error('Session expiree');
      } else if (err.status !== 0) {
        toast.error(err.error?.message || `Erreur ${err.status}`);
      }
      return throwError(() => err);
    })
  );
};
\end{lstlisting}

\section{AuthGuard}

\begin{lstlisting}[style=tsstyle]
export const authGuard: CanActivateFn = () => {
  const auth = inject(AuthService);
  const router = inject(Router);
  if (auth.isAuthenticated) return true;
  router.navigate(['/login']);
  return false;
};
\end{lstlisting}

Trois routes sont protégées par cette garde : \texttt{/cart},
\texttt{/orders} et \texttt{/profile}. Le panier d'un visiteur non
connecté n'est donc jamais persisté côté serveur.

\section{Aperçu de l'interface}

L'interface adopte les codes du e-commerce moderne : barre de navigation
fixe en tête, grille de cartes produits responsive (1 / 2 / 4 colonnes
selon la largeur d'écran), badges colorés indiquant la disponibilité
stock. Bootstrap 5 fournit l'ossature CSS et les composants standards
(modales, alertes, formulaires).

% =====================================================================
%  CHAPITRE 7 : TESTS ET VALIDATION
% =====================================================================
\chapter{Tests et validation}

\section{Stratégie de tests}

Suivant la \textbf{pyramide de tests} de Mike Cohn, nous avons concentré
l'effort sur la base de la pyramide : les tests unitaires des services
métier. Ces tests sont rapides, déterministes et faciles à maintenir.

\section{Outils utilisés}

\begin{itemize}
  \item \textbf{JUnit 5} (Jupiter) : framework de test
  \item \textbf{Mockito 5} : isolation des dépendances
  \item \textbf{AssertJ} : assertions fluentes (\texttt{assertThat...})
  \item \textbf{Spring Boot Test} : pour les tests d'intégration éventuels
\end{itemize}

\section{Suites de tests implémentées}

\begin{table}[htbp]
\centering
\caption{Tests unitaires E-Store}
\label{tab:tests}
\small
\begin{tabularx}{\textwidth}{lccX}
\toprule
\textbf{Suite} & \textbf{Tests} & \textbf{Statut} & \textbf{Couverture}\\
\midrule
\texttt{ProductServiceTest} & 4 & PASS & findById (présent / absent), recherche (avec / sans mot-clé)\\
\texttt{CartServiceTest}    & 2 & PASS & Ajout avec stock suffisant et insuffisant\\
\texttt{OrderServiceTest}   & 3 & PASS & Panier vide, validation complète, historique\\
\midrule
\textbf{Total}              & \textbf{9} & \textbf{0 failure} & BUILD SUCCESS\\
\bottomrule
\end{tabularx}
\end{table}

\section{Exemple de test}

\begin{lstlisting}[style=javastyle, caption=Test du checkout transactionnel]
@ExtendWith(MockitoExtension.class)
@MockitoSettings(strictness = Strictness.LENIENT)
class OrderServiceTest {

    @Mock private OrderRepository orderRepository;
    @Mock private CartService cartService;
    @Mock private InventoryService inventoryService;
    @InjectMocks private OrderService orderService;

    @Test
    void checkout_panierValide_creeCommandeEtViderPanier() {
        // Given
        cart.getItems().add(item);
        when(orderRepository.save(any(Order.class)))
            .thenAnswer(inv -> { Order o = inv.getArgument(0);
                                  o.setId(500L); return o; });
        // When
        OrderDto dto = orderService.checkout();
        // Then
        assertThat(dto.getId()).isEqualTo(500L);
        assertThat(dto.getStatus()).isEqualTo(OrderStatus.CONFIRMED);
        assertThat(dto.getTotalAmount())
            .isEqualByComparingTo(new BigDecimal("25000.00"));
        verify(inventoryService).decrement(10L, 2);
        verify(cartService).clearCart(cart);
    }

    @Test
    void checkout_panierVide_lanceException() {
        assertThatThrownBy(() -> orderService.checkout())
            .isInstanceOf(BusinessException.class)
            .hasMessageContaining("vide");
    }
}
\end{lstlisting}

\section{Résultats d'exécution}

\begin{lstlisting}[style=shellstyle]
$ ./mvnw test
[INFO] Tests run: 3, Failures: 0, Errors: 0 -- OrderServiceTest
[INFO] Tests run: 4, Failures: 0, Errors: 0 -- ProductServiceTest
[INFO] Tests run: 2, Failures: 0, Errors: 0 -- CartServiceTest
[INFO] Tests run: 9, Failures: 0, Errors: 0, Skipped: 0
[INFO] BUILD SUCCESS
\end{lstlisting}

% =====================================================================
%  CHAPITRE 8 : DÉPLOIEMENT ET DEVOPS
% =====================================================================
\chapter{Déploiement et DevOps}

\section{Profils Spring}

L'application fournit deux profils d'exécution :

\begin{itemize}
  \item \textbf{dev} (par défaut) : H2 in-memory, MongoDB optionnel.
    Démarrage instantané sans installation préalable.
  \item \textbf{prod} : MySQL 8 + MongoDB 7, accessibles via Docker Compose.
\end{itemize}

\section{Conteneurisation Docker}

Le fichier \texttt{docker-compose.yml} orchestre quatre services :

\begin{lstlisting}[style=propertiesstyle, caption=docker-compose.yml (extrait)]
services:
  mysql:
    image: mysql:8
    environment:
      MYSQL_ROOT_PASSWORD: root
      MYSQL_DATABASE: estore
    ports: ["3306:3306"]
    volumes: [mysql-data:/var/lib/mysql]

  mongo:
    image: mongo:7
    ports: ["27017:27017"]
    volumes: [mongo-data:/data/db]

  phpmyadmin:
    image: phpmyadmin:latest
    ports: ["8081:80"]
    depends_on: [mysql]

  mongo-express:
    image: mongo-express:latest
    ports: ["8082:8081"]
    depends_on: [mongo]
\end{lstlisting}

Une commande suffit alors à provisionner l'environnement complet :

\begin{lstlisting}[style=shellstyle]
$ docker compose up -d
\end{lstlisting}

\section{Démarrage en trois commandes}

\begin{lstlisting}[style=shellstyle]
# 1. Bases de donnees (optionnel : profil prod uniquement)
docker compose up -d

# 2. Backend
cd estore-backend && ./mvnw spring-boot:run

# 3. Frontend
cd estore-frontend && npm install && npm start
\end{lstlisting}

\section{Versionnement Git et publication GitHub}

L'historique du projet est versionné par Git, organisé en
\textbf{19 commits} sémantiques (\texttt{feat}, \texttt{fix},
\texttt{docs}, \texttt{chore}). Le projet est publié sur GitHub à l'adresse
\url{https://github.com/akrambelmoussa-etu-byte/e-store}.

% =====================================================================
%  CHAPITRE 9 : DIFFICULTÉS RENCONTRÉES ET SOLUTIONS
% =====================================================================
\chapter{Difficultés rencontrées et solutions}

Le développement n'a pas été un long fleuve tranquille. Plusieurs
obstacles techniques significatifs ont été rencontrés ; les surmonter a
été l'occasion d'approfondir notre compréhension de l'écosystème.

\section{Difficulté n°1 : strictness Mockito}

\subsection{Symptôme}
Lors du premier lancement de \texttt{mvn test}, deux tests échouent avec
le message :
\begin{quote}
\texttt{UnnecessaryStubbingException: Unnecessary stubbings detected.}
\end{quote}

\subsection{Cause}
Mockito 5 applique par défaut le mode strict. Toute instruction
\texttt{when(...).thenReturn(...)} déclarée mais non vérifiée durant
l'exécution du test provoque une erreur. Or, nos méthodes
\texttt{@BeforeEach} configuraient des stubs partagés entre plusieurs tests.

\subsection{Solution}
Annoter chaque classe de test avec
\texttt{@MockitoSettings(strictness = Strictness.LENIENT)}, ce qui
relâche la vérification tout en conservant la sécurité d'écriture.

\begin{lstlisting}[style=javastyle]
@ExtendWith(MockitoExtension.class)
@MockitoSettings(strictness = Strictness.LENIENT)
class CartServiceTest { ... }
\end{lstlisting}

\section{Difficulté n°2 : LazyInitializationException}

\subsection{Symptôme}
Lors du retour d'une entité \texttt{User} via le contrôleur, Hibernate
levait une \texttt{LazyInitializationException} sur le \texttt{Profile}
chargé en mode \texttt{LAZY}.

\subsection{Cause}
La sérialisation Jackson tentait de naviguer dans la relation après la
fermeture de la session JPA.

\subsection{Solution}
Adoption systématique du pattern \textbf{DTO} : chaque endpoint renvoie
un objet de transfert (\texttt{UserDetailsDto}), construit dans la couche
service alors que la session est encore ouverte. Avantage collatéral :
les DTOs masquent les champs sensibles (password) et stabilisent l'API.

\section{Difficulté n°3 : installation de Docker Desktop}

\subsection{Symptôme}
Lors de l'installation de Docker Desktop sous Windows 11, l'erreur
suivante bloque le processus :
\begin{quote}\itshape
For security reasons C:\textbackslash ProgramData\textbackslash DockerDesktop must be owned
by an elevated account.
\end{quote}

\subsection{Cause}
Un dossier résiduel d'une installation antérieure subsistait avec des
permissions incorrectes.

\subsection{Solutions implémentées}

\textbf{Solution A — Brute :} suppression manuelle en mode administrateur :
\begin{lstlisting}[style=shellstyle]
rmdir /s /q "C:\ProgramData\DockerDesktop"
rmdir /s /q "C:\ProgramData\Docker"
\end{lstlisting}
puis réinstallation.

\textbf{Solution B — Élégante :} mise en place d'un \textbf{profil dev}
basé sur H2 in-memory, qui permet de démarrer l'application sans aucune
installation préalable. Cette dualité Docker/H2 a été pensée dès la
conception et a fait ses preuves le jour de la démonstration.

\begin{tipbox}
Cette difficulté a été convertie en force pédagogique : nous démontrons
au jury notre capacité à anticiper les problèmes d'environnement et à
prévoir des chemins de repli (\textit{fallback}) — une compétence
clairement attendue en milieu professionnel.
\end{tipbox}

\section{Difficulté n°4 : encodage des tokens JWT}

\subsection{Symptôme}
Lors de la génération du premier JWT, la bibliothèque \texttt{jjwt} levait
une \texttt{WeakKeyException} : la clé secrète configurée comportait moins
de 256 bits.

\subsection{Solution}
Implémenter une logique de fallback : tenter d'abord un décodage Base64
de la clé, et en cas d'échec ou de longueur insuffisante, utiliser
l'encodage UTF-8 brut.

\begin{lstlisting}[style=javastyle]
byte[] bytes;
try {
    bytes = Decoders.BASE64.decode(secret);
    if (bytes.length < 32) throw new IllegalArgumentException();
} catch (Exception e) {
    bytes = secret.getBytes(StandardCharsets.UTF_8);
}
this.key = Keys.hmacShaKeyFor(bytes);
\end{lstlisting}

% =====================================================================
%  CHAPITRE 10 : RÉSULTATS
% =====================================================================
\chapter{Résultats et démonstration}

\section{Bilan fonctionnel}

L'application livrée couvre intégralement le périmètre du cahier des
charges. Le tableau~\ref{tab:bilan} en synthétise les fonctionnalités.

\begin{table}[htbp]
\centering
\caption{Fonctionnalités livrées}
\label{tab:bilan}
\small
\begin{tabularx}{\textwidth}{lXc}
\toprule
\textbf{Fonctionnalité} & \textbf{Description} & \textbf{Statut}\\
\midrule
Inscription / Connexion & Email + mot de passe, JWT 24h            & Livré \\
Profil utilisateur      & Consultation et modification             & Livré \\
Catalogue paginé        & Recherche, filtre par catégorie          & Livré \\
Fiche produit           & Détail, image, stock, avis               & Livré \\
Panier persistant       & Ajout, modification, vidage              & Livré \\
Validation commande     & Transactionnelle ACID                    & Livré \\
Historique commandes    & Consultation détaillée                   & Livré \\
Avis produits           & Dépôt et consultation (MongoDB)          & Livré \\
Administration          & CRUD catégories / produits / stock       & Livré \\
\bottomrule
\end{tabularx}
\end{table}

\section{Métriques de qualité}

\begin{table}[htbp]
\centering
\caption{Métriques quantitatives du projet}
\label{tab:metrics}
\small
\begin{tabularx}{\textwidth}{Xr}
\toprule
\textbf{Indicateur} & \textbf{Valeur}\\
\midrule
Lignes de code Java                              & $\sim$ 3 500 \\
Lignes de code TypeScript                        & $\sim$ 2 000 \\
Nombre de classes Java                           & 71 \\
Nombre de composants Angular                     & 14 \\
Nombre d'endpoints REST                          & 24 \\
Tests unitaires                                  & 9 (100\% passants) \\
Bundle JS initial                                & 13.84 ko \\
Temps de démarrage backend (H2)                  & $\approx$ 4 s \\
Temps de réponse API (médiane, dev)              & $\approx$ 30 ms \\
Commits Git                                      & 19 \\
\bottomrule
\end{tabularx}
\end{table}

\section{Captures d'écran de la démonstration}

Les captures d'écran de l'application en fonctionnement sont disponibles
dans le dossier \texttt{docs/screenshots/} du dépôt GitHub. Le scénario
de démonstration suit huit étapes :

\begin{enumerate}
  \item Page catalogue avec 12 produits seed
  \item Filtrage par catégorie « Sport »
  \item Fiche détaillée d'un produit
  \item Tentative d'ajout au panier non authentifié $\to$ redirection
  \item Connexion avec le compte de test
  \item Ajout de deux produits au panier
  \item Validation de la commande, redirection vers /orders
  \item Dépôt d'un avis 5 étoiles, persistance MongoDB
\end{enumerate}

% =====================================================================
%  CONCLUSION GÉNÉRALE
% =====================================================================
\chapter*{Conclusion générale}
\addcontentsline{toc}{chapter}{Conclusion générale}
\markboth{CONCLUSION GÉNÉRALE}{}

\section*{Bilan}

Au terme de ce projet, nous avons mené à bien la conception et la
réalisation d'une application e-commerce \textbf{intégralement
fonctionnelle}, couvrant le parcours utilisateur du visiteur anonyme
jusqu'à la commande validée. L'architecture en \textbf{trois couches × cinq
domaines} a été rigoureusement respectée, fournissant une base saine pour
toute évolution ultérieure.

Sur le plan technique, les objectifs initiaux ont tous été atteints :
\begin{itemize}
  \item Architecture full-stack moderne avec Spring Boot 3.3 et Angular 17.
  \item Persistance hybride MySQL + MongoDB opérationnelle.
  \item Sécurité robuste par BCrypt + JWT + \texttt{@PreAuthorize}.
  \item Suite de tests JUnit 5 / Mockito en BUILD SUCCESS.
  \item Conteneurisation Docker Compose, profil de fallback H2.
\end{itemize}

Sur le plan humain, le travail en binôme nous a permis d'expérimenter
\textit{in vivo} les pratiques du développement collaboratif :
versionnement Git, conventions de commit, revue de code, répartition des
tâches, communication asynchrone. Ces compétences relationnelles, souvent
sous-estimées, sont déterminantes en milieu professionnel.

\section*{Apports personnels}

Au-delà de la technique, ce projet nous a confronté à la \textbf{réalité
de l'incertitude} qui caractérise tout développement logiciel : choix
d'architecture aux conséquences durables, débogage de problèmes inattendus
(la mésaventure Docker en est l'archétype), arbitrages entre élégance du
code et délais. Apprendre à composer avec cette incertitude, plutôt qu'à
la fuir, est probablement le bénéfice le plus durable de cette expérience.

\section*{Perspectives}

Le projet livré n'est pas un produit fini mais un \textbf{socle évolutif}.
Plusieurs pistes d'extension ont été identifiées :

\begin{itemize}
  \item \textbf{Paiement en ligne} : intégration de Stripe ou des
    passerelles marocaines (CMI, Maroc Telecommerce).
  \item \textbf{Notifications transactionnelles} : envoi automatique
    d'emails (Spring Mail) ou de SMS (Twilio) lors de la confirmation de
    commande.
  \item \textbf{Recommandations} : suggestion produits basée sur
    l'historique d'achat (collaborative filtering simple, ou Spring AI).
  \item \textbf{Internationalisation} : passage du français à l'arabe ou
    l'anglais via Angular i18n.
  \item \textbf{Déploiement cloud} : conteneurisation et déploiement sur
    Heroku, Render ou AWS Free Tier.
  \item \textbf{Application mobile} : portage Ionic ou Flutter consommant
    la même API.
  \item \textbf{Observabilité} : intégration Prometheus + Grafana pour les
    métriques applicatives.
\end{itemize}

Plus fondamentalement, l'architecture en couches et en domaines que nous
avons mise en place rend ces évolutions \textbf{additives} : il sera
possible de les introduire sans casser l'existant. C'est, à nos yeux, la
preuve la plus convaincante que les choix architecturaux de départ étaient
les bons.

% =====================================================================
%  BIBLIOGRAPHIE
% =====================================================================
\chapter*{Bibliographie et webographie}
\addcontentsline{toc}{chapter}{Bibliographie et webographie}
\markboth{BIBLIOGRAPHIE}{}

\subsection*{Ouvrages}

\begin{enumerate}
  \item Craig Walls, \textit{Spring Boot in Action}, 2nd edition,
    Manning Publications, 2024.
  \item Robert C. Martin, \textit{Clean Code: A Handbook of Agile Software
    Craftsmanship}, Prentice Hall, 2008.
  \item Eric Evans, \textit{Domain-Driven Design: Tackling Complexity in
    the Heart of Software}, Addison-Wesley, 2003.
  \item Martin Fowler, \textit{Patterns of Enterprise Application
    Architecture}, Addison-Wesley, 2002.
\end{enumerate}

\subsection*{Documentations officielles}

\begin{enumerate}
  \item Spring Framework Reference Documentation, \url{https://docs.spring.io/spring-framework/reference/}
  \item Spring Boot Reference Documentation, \url{https://docs.spring.io/spring-boot/}
  \item Spring Data JPA, \url{https://docs.spring.io/spring-data/jpa/reference/}
  \item Spring Data MongoDB, \url{https://docs.spring.io/spring-data/mongodb/reference/}
  \item Angular Documentation, \url{https://angular.dev/}
  \item MongoDB Manual, \url{https://www.mongodb.com/docs/manual/}
  \item RFC 7519 — JSON Web Token (JWT), \url{https://www.rfc-editor.org/rfc/rfc7519}
  \item OWASP Authentication Cheat Sheet, \url{https://cheatsheetseries.owasp.org/cheatsheets/Authentication_Cheat_Sheet.html}
\end{enumerate}

\subsection*{Articles et tutoriels}

\begin{enumerate}
  \item Baeldung, ``A Guide to JPA with Spring'', \url{https://www.baeldung.com/the-persistence-layer-with-spring-and-jpa}
  \item Baeldung, ``Spring Security with JWT'', \url{https://www.baeldung.com/spring-security-oauth-jwt}
  \item Angular Blog, ``Introducing Angular Signals'', \url{https://blog.angular.io/}
\end{enumerate}

% =====================================================================
%  ANNEXES
% =====================================================================
\appendix

\chapter{Configuration \texttt{application.properties}}

\begin{lstlisting}[style=propertiesstyle]
# === application.properties (commun) ===
spring.application.name=estore-backend
spring.profiles.active=dev
server.port=8080

# JPA
spring.jpa.hibernate.ddl-auto=update
spring.jpa.open-in-view=false

# MongoDB
spring.data.mongodb.uri=mongodb://localhost:27017/estore

# JWT
jwt.secret=change-me-with-a-256-bits-secret-key-for-production
jwt.expiration=86400000

# === application-dev.properties (H2 in-memory) ===
spring.datasource.url=jdbc:h2:mem:estore;MODE=MySQL
spring.datasource.driver-class-name=org.h2.Driver
spring.jpa.database-platform=org.hibernate.dialect.H2Dialect
spring.h2.console.enabled=true

# === application-prod.properties (MySQL) ===
spring.datasource.url=jdbc:mysql://localhost:3306/estore?createDatabaseIfNotExist=true
spring.datasource.username=root
spring.datasource.password=root
spring.jpa.database-platform=org.hibernate.dialect.MySQLDialect
\end{lstlisting}

\chapter{Liste exhaustive des endpoints REST}

\begin{longtable}{p{1.2cm}p{4.5cm}p{6cm}c}
\caption{Liste complète des endpoints REST}\label{tab:endpoints-full}\\
\toprule
\textbf{Verbe} & \textbf{URL} & \textbf{Description} & \textbf{Auth.}\\
\midrule
\endfirsthead
\toprule
\textbf{Verbe} & \textbf{URL} & \textbf{Description} & \textbf{Auth.}\\
\midrule
\endhead
POST   & /api/auth/register     & Inscription                  & Public \\
POST   & /api/auth/login        & Connexion (renvoie JWT)      & Public \\
GET    & /api/users/me          & Profil courant               & USER \\
PUT    & /api/users/me          & Mise à jour profil           & USER \\
GET    & /api/categories        & Liste catégories             & Public \\
GET    & /api/categories/\{id\} & Détail catégorie             & Public \\
POST   & /api/categories        & Création                     & ADMIN \\
PUT    & /api/categories/\{id\} & Modification                 & ADMIN \\
DELETE & /api/categories/\{id\} & Suppression                  & ADMIN \\
GET    & /api/products          & Catalogue + filtres + pagination & Public \\
GET    & /api/products/\{id\}   & Fiche produit                & Public \\
POST   & /api/products          & Création                     & ADMIN \\
PUT    & /api/products/\{id\}   & Modification                 & ADMIN \\
DELETE & /api/products/\{id\}   & Suppression                  & ADMIN \\
GET    & /api/inventory/\{id\}  & Lecture stock                & Public \\
PUT    & /api/inventory/\{id\}  & Mise à jour stock            & ADMIN \\
GET    & /api/cart              & Panier courant               & USER \\
POST   & /api/cart/add          & Ajout                        & USER \\
PUT    & /api/cart/update       & Modification quantité        & USER \\
DELETE & /api/cart/remove/\{id\}& Retrait article              & USER \\
DELETE & /api/cart/clear        & Vidage                       & USER \\
POST   & /api/orders            & Validation commande          & USER \\
GET    & /api/orders            & Historique                   & USER \\
GET    & /api/orders/\{id\}     & Détail commande              & USER \\
POST   & /api/reviews           & Dépôt avis                   & USER \\
GET    & /api/reviews/product/\{id\} & Avis d'un produit        & Public \\
GET    & /api/reviews/user/\{id\} & Avis d'un utilisateur       & Public \\
\bottomrule
\end{longtable}

\chapter{Comptes de test}

Le \texttt{DataSeeder} crée automatiquement deux comptes au premier
démarrage si la base est vide :

\begin{table}[htbp]
\centering
\small
\begin{tabular}{lll}
\toprule
\textbf{Rôle} & \textbf{Email} & \textbf{Mot de passe}\\
\midrule
ADMIN & admin@estore.ma & Admin@123 \\
USER  & user@estore.ma  & User@123  \\
\bottomrule
\end{tabular}
\end{table}

\chapter{Procédure de démarrage rapide}

\begin{lstlisting}[style=shellstyle]
# 1. Cloner le depot
git clone https://github.com/akrambelmoussa-etu-byte/e-store.git
cd e-store

# 2a. Demarrer en mode dev (H2 + Mongo optionnel)
cd estore-backend
./mvnw spring-boot:run

# 2b. (Alternative) Demarrer en mode prod
docker compose up -d
cd estore-backend
./mvnw spring-boot:run -Dspring.profiles.active=prod

# 3. Frontend (terminal separe)
cd estore-frontend
npm install
npm start

# 4. Ouvrir http://localhost:4200
\end{lstlisting}

\end{document}
